\documentclass{article}

\PassOptionsToPackage{numbers, sort, compress}{natbib}
\usepackage[main,final]{neurips_2025}

\usepackage{cmap}                       % корректный поиск и копирование кириллицы из PDF
\usepackage[T2A,T1]{fontenc}
\usepackage[utf8]{inputenc}
\usepackage[english,russian]{babel}
\usepackage[unicode]{hyperref}
\usepackage{url}
\usepackage{booktabs}
\usepackage{nicefrac}
\usepackage{microtype}
\usepackage{xcolor}

%%%

\usepackage{subcaption}
\usepackage{graphicx}
\usepackage{multirow}
\usepackage{amsmath,amssymb,amsfonts}
\usepackage{amsthm}
\usepackage{mathrsfs}
\usepackage{xcolor}
\usepackage{textcomp}
\usepackage{manyfoot}
\usepackage{algorithm}
\usepackage{algorithmicx}
\usepackage{algpseudocode}
\usepackage{listings}

% Шрифт Times (ptm), задаваемый стилем NeurIPS, не содержит кириллицы в стандартной
% поставке TeX Live. При наличии пакета tempora используется Times-подобный шрифт
% с кириллицей, иначе -- Computer Modern (cm-super), содержащий кодировку T2A.
\IfFileExists{tempora.sty}{\usepackage{tempora}}{\renewcommand{\rmdefault}{cmr}}

% Подпись внизу первой страницы вместо стандартной надписи NeurIPS.
\makeatletter
\renewcommand{\@noticestring}{%
  Выпускная квалификационная работа бакалавра. Факультет вычислительной математики
  и кибернетики МГУ имени М.\,В.~Ломоносова. Оформление по шаблону курса
  <<Моя первая научная статья>>. Черновик, 2026.%
}
% Заголовок аннотации по-русски (стиль NeurIPS задаёт его жёстко).
\renewenvironment{abstract}%
{%
  \vskip 0.075in%
  \centerline%
  {\large\bf Аннотация}%
  \vspace{0.5ex}%
  \begin{quote}%
}
{
  \par%
  \end{quote}%
  \vskip 1ex%
}
\makeatother

\newtheorem{theorem}{Теорема} % continuous numbers
%%\newtheorem{theorem}{Теорема}[section] % sectionwise numbers
%% optional argument [theorem] produces theorem numbering sequence instead of independent numbers for Proposition
\newtheorem{proposition}[theorem]{Предложение}%
\newtheorem{lemma}{Лемма}%
%%\newtheorem{proposition}{Предложение} % to get separate numbers for theorem and proposition etc.

\newtheorem{example}{Пример}
\newtheorem{remark}{Замечание}

\newtheorem{definition}{Определение}
\newtheorem{assumption}{Предположение}

% Обозначения, используемые в тексте буквально (см. docs/problem_statement.md).
\newcommand{\I}{\mathbb{I}}                         % индикатор
\newcommand{\Inv}{\mathrm{Inv}}
\newcommand{\Heval}{\mathcal{H}_{\mathrm{eval}}}
\newcommand{\fprod}{f_{\mathrm{prod}}}
\newcommand{\NA}{\mathrm{NA}}
\newcommand{\norm}{\mathrm{norm}}
\newcommand{\EM}{\mathrm{EM}}
\newcommand{\CP}{\mathrm{CP}}
\newcommand{\CR}{\mathrm{CR}}
\newcommand{\AR}{\mathrm{AR}}
\newcommand{\FAR}{\mathrm{FAR}}
\newcommand{\mcorr}{m_J^{\mathrm{corr}}}

\lstdefinelanguage{json}{
  basicstyle=\ttfamily\footnotesize,
  showstringspaces=false,
  breaklines=true,
  frame=single,
  string=[s]{"}{"},
  stringstyle=\color{blue!60!black},
  comment=[l]{//},
  commentstyle=\color{gray},
}

%%%


\title{Система оценки качества ответов больших языковых моделей по новостным текстам стран Восточной Азии\\
  \vspace{0.5ex}{\large A System for Evaluating the Quality of Large Language Model Answers on East Asian News Texts}}


% The \author macro works with any number of authors. There are two commands
% used to separate the names and addresses of multiple authors: \And and \AND.
%
% Using \And between authors leaves it to LaTeX to determine where to break the
% lines. Using \AND forces a line break at that point. So, if LaTeX puts 3 of 4
% authors names on the first line, and the last on the second line, try using
% \AND instead of \And before the third author name.


\author{
  Багров Александр Михайлович\\
  Факультет вычислительной математики и кибернетики\\
  МГУ имени М.\,В.~Ломоносова\\
  Москва, Россия\\
  \texttt{alexbagrov04@gmail.com}\\
  \And
  Научный руководитель: Воронцов Константин Вячеславович, д.\,ф.-м.\,н.\\
  Факультет вычислительной математики и кибернетики\\
  МГУ имени М.\,В.~Ломоносова\\
  Москва, Россия\\
}


\begin{document}


\maketitle

\begin{abstract}
Промышленные вопросно-ответные системы, дополняющие большую языковую модель поиском по новостному корпусу, внедряются быстрее, чем появляются средства количественной проверки их ответов: существующие наборы вопросов либо англоязычны, либо построены из вопросов с меняющимся во времени ответом, а согласие судьи-модели с людьми и согласие самих аннотаторов в них, как правило, не сообщаются. В работе строится система оценки качества ответов больших языковых моделей на вопросы по новостным текстам стран Восточной Азии в кросс-языковой постановке: вопрос задаётся на русском языке, а источники~--- китайскоязычные СМИ. Объектом измерения является ответ $\hat a = f(q)$ на вопрос $q$ о зафиксированном снимке корпуса, а не порождающая его модель. Система оценки включает: (1)~формальную постановку с предикатом инвариантности $\Inv(q)$, ограничивающим выборку статическими вопросами, ответ на которые не меняется во времени; (2)~эталонную выборку $G = G^+ \sqcup G^{\varnothing}$ (не менее 150 вопросов, включая вопросы с правильным ответом-отказом) с эталонными ответами, алиасами, опорными документами, независимой разметкой тремя аннотаторами и мерой согласия Криппендорфа $\alpha \ge 0{,}67$; (3)~семейство метрик $M$ по эталону, цитирования, отказов и поиска, включая судью на основе языковой модели, который проходит мета-оценку раздельно на надёжность (воспроизводимость оценок) и валидность (корреляция Спирмена с человеческими оценками не ниже $0{,}5$), после чего отбирается валидированное подсемейство $M'$; (4)~статистический протокол сравнения систем по композитному показателю $Q(f)$~--- парный бутстреп, критерий Уилкоксона с поправкой Холма и $\delta$~Клиффа. Объект исследования~--- промышленная система <<Chinese Media Analysis>> (ИВ~РАН) над корпусом СМИ КНР, Тайваня, Гонконга и Сингапура (23,9~млн фрагментов); она сравнивается с шестью открытыми моделями в режимах без поиска и с общим ретривером (13~систем). Расширение на японские и корейские источники~--- предмет дальнейшей работы. Эталонная выборка, конфигурации систем, сырые генерации и код публикуются.
\end{abstract}

\textbf{Ключевые слова:} большие языковые модели, генерация с дополнением извлечёнными данными (RAG), оценка качества ответов, эталонная выборка, статические вопросы, кросс-языковой поиск, китайскоязычные СМИ, судья на основе языковой модели, мета-оценка судьи, самопредпочтение, согласие аннотаторов.

\section{Введение}\label{sec:intro}

Новостные тексты стран Восточной Азии~--- основной источник сведений для востоковедов, аналитиков и журналистов, однако объём потока и языковой барьер делают ручной мониторинг невозможным. Естественным инструментом стали системы генерации с дополнением извлечёнными данными (RAG)~\cite{lewis2020rag}: поисковый компонент находит фрагменты новостей, а большая языковая модель (LLM) формирует ответ на естественном языке с цитированием источников. Один из таких инструментов~--- промышленная система <<Chinese Media Analysis>>, разработанная в Институте востоковедения РАН и описанная в работе~\cite{bagrov2026digital}; она работает над корпусом китайскоязычных СМИ КНР, Тайваня, Гонконга и Сингапура объёмом 23,9~млн фрагментов и отвечает на русскоязычные вопросы.

Качество ответов таких систем до сих пор оценивалось качественно. Транскрипты сеансов показывают характерные дефекты: ложную атрибуцию (факты 2021--2025 годов со ссылкой на статью 2019 года), номера цитат, отсутствующие в списке источников, нерелевантные фрагменты с высокой косинусной близостью, непереведённые китаизмы. Без измеримых критериев невозможно ответить на вопросы <<стала ли система лучше после смены модели>> и <<оправдан ли поиск по корпусу по сравнению с ответом модели без поиска>>. Существующие средства оценки RAG~\cite{es2023ragas,saadfalcon2023ares} опираются на LLM-судью без проверки его согласия с человеком на русскоязычных ответах по китайскому корпусу, хотя надёжность и валидность судьи~--- разные свойства, и обе требуют измерения~\cite{norman2026reliability,yamauchi2025empirical}, а согласие судьи с людьми на русском языке заметно ниже, чем на английском~\cite{pugachev2025repa}. Доступные наборы вопросов~\cite{chen2023rgb,lyu2024crudrag,chen2025chronoqa} либо составлены на английском и китайском языках, либо содержат зависящие от времени вопросы, из-за чего эталон устаревает~\cite{piryani2026recency}; вопросы, правильный ответ на которые~--- отказ, в них, как правило, отсутствуют, хотя именно готовность отвечать без опоры на корпус отличает системы сильнее, чем точность~\cite{dorosario2026penalty,kirichenko2025abstentionbench}.

В настоящей работе строится система оценки, объектом измерения в которой является \emph{ответ} $\hat a=f(q)$ на вопрос $q$ о новостном корпусе, а не порождающая его модель; сравнение систем~--- следствие применения системы оценки. Рассматриваются только статические вопросы, ответ на которые не меняется во времени (класс never-changing в таксономии~\cite{vu2023freshllms}), что делает эталонную выборку долговечной. Программная система~\cite{bagrov2026digital} является объектом исследования и не входит в результаты работы. Имеющийся в её репозитории прототип модуля оценки (LLM-судья по семи измерениям, метрики поиска), не запускавшийся и не валидированный, рассматривается как кандидатное подсемейство метрик, подлежащее валидации.

\textbf{Вклад работы.} Основные результаты состоят в следующем:
\begin{itemize}
    \item формальная постановка задачи оценки качества ответов на статические вопросы по новостному корпусу: предикат инвариантности $\Inv(q)$, класс гипотез $\mathcal{H}$, семейство метрик $M$, композитный показатель $Q(f)$ и критерий сравнения со статистической проверкой;
    \item эталонная выборка $G=G^+\sqcup G^{\varnothing}$ статических вопросов с каноническими ответами, алиасами, опорными документами и вопросами, требующими отказа; протокол независимой разметки $R=3$ аннотаторами с мерой согласия Криппендорфа $\alpha$ и адъюдикацией;
    \item валидированное семейство метрик $M'$: метрика правильности по эталону $\mcorr$, детерминированные метрики $\EM$, F1, BERTScore, метрики цитирования $\CP_{\mathrm{id}}$, $\CR_{\mathrm{id}}$, $\CP_{\mathrm{ent}}$ и метрики отказов $\AR$, $1-\FAR$; отбор метрик по коэффициенту Спирмена $\rho_S$ с человеческими оценками;
    \item первое систематическое сравнение промышленной RAG-системы с открытыми моделями в режимах без поиска и с общим ретривером с парным бутстрепом, критерием Уилкоксона и поправкой Холма;
    \item открытые данные и код: эталонная выборка, конфигурации систем, сырые генерации, снимок списка моделей провайдера и скрипты вычисления метрик и статистики.
\end{itemize}

\textbf{Структура работы.} В разделе~\ref{sec:rw} рассматриваются смежные работы. Раздел~\ref{sec:prelim} вводит обозначения, предположения и формальную постановку задачи. В разделе~\ref{sec:method} описана методика построения системы оценки, в разделе~\ref{sec:exp}~--- план экспериментов и гипотезы. Раздел~\ref{sec:disc} посвящён рискам и ограничениям, раздел~\ref{sec:concl} содержит выводы. В приложении~\ref{app} приведены правила исключения вопросов и схема записи эталонной выборки.

\section{Обзор литературы}\label{sec:rw}

\textbf{Оценка систем генерации с дополнением извлечёнными данными.}
Фреймворк RAGAS~\cite{es2023ragas} ввёл безэталонные метрики faithfulness, answer relevancy и context relevance, вычисляемые LLM-судьёй; ARES~\cite{saadfalcon2023ares} дообучает лёгкие судьи на синтетических данных и калибрует их по небольшой размеченной выборке методом prediction-powered inference. Бенчмарк RGB~\cite{chen2023rgb} проверяет устойчивость к шуму, отказ при отсутствии информации и интеграцию нескольких источников на английских и китайских новостях; CRUD-RAG~\cite{lyu2024crudrag} строит задачи над китайскими новостями по операциям create, read, update, delete. Работы по треку TREC 2024 RAG показали, что судья надёжно оценивает поддержку утверждений цитатами~\cite{thakur2025assessing} и что оценка по атомарным фактам эталона (nugget) воспроизводит ручную разметку~\cite{pradeep2025nugget}, однако правильность цитирования не гарантирует, что модель действительно опиралась на источник~\cite{wallat2025correctness}, а вводящая в заблуждение выдача может сделать ответ с поиском хуже ответа без него~\cite{zeng2025raguard}. Метрики выдачи ретривера согласуются с полнотой ответа лишь при совпадении целей поиска и генерации~\cite{samuel2026beyond}. Эти работы задают словарь метрик, но не содержат проверки согласия судьи с человеком для русскоязычных ответов и не выделяют статические вопросы. Обзор архитектур RAG и мотивация абляции <<с поиском/без поиска>> восходят к~\cite{lewis2020rag}; на редких фактах поиск даёт наибольший выигрыш, тогда как на популярных модель без поиска конкурентоспособна~\cite{mallen2022popqa}.

\textbf{Большие языковые модели как судьи.}
Работа~\cite{zheng2023judging} показала, что сильные модели согласуются с человеческими предпочтениями на уровне попарного согласия людей, но подвержены позиционному смещению, смещению к длинному ответу и ограниченной способности к проверке фактов; в чистом сценарии с высоким согласием людей даже лучшие судьи остаются заметно ниже межчеловеческого согласия~\cite{thakur2024judging}. В~\cite{panickssery2024selfpreference} установлено самопредпочтение: модель выше оценивает тексты, которые распознаёт как свои; позднее показано, что смещение распространяется на родственные модели того же семейства~\cite{li2025preference} и сохраняется даже при бинарной рубрике с объективными критериями~\cite{pombal2026selfpreference}, а панель судей разных семейств уменьшает его~\cite{verga2024replacing}. Крупнейшая мета-оценка судей~\cite{norman2026reliability} различает надёжность (воспроизводимость) и валидность (согласие с людьми с поправкой на случайность) и показывает, что доля совпадений завышает согласие на десятки процентных пунктов относительно $\kappa$. При оценке правильности ответа относительно эталона судья коррелирует с людьми до $0{,}85$ против $0{,}22$ у точного совпадения и $0{,}40$ у F1~\cite{ho2025reassessing}, а половина из 54 судей воспроизводит человеческое суждение о правильности ответов RAG~\cite{han2025verdict}; для мета-оценки судей цитирующих систем выпущен ресурс Auto-Judge~\cite{farzi2026autojudge}. Замену человека судьёй можно обосновать статистически по небольшой размеченной подвыборке~\cite{calderon2025alternative}. Согласие судьи, установленное на английском языке, не переносится на другие языки~\cite{fu2025multilingual}, и для русского оно заметно ниже~\cite{pugachev2025repa}; мета-оценочные наборы для неанглийского RAG появились лишь недавно~\cite{cruzblandon2025memerag}. Для настоящей работы это означает, что судья не должен совпадать ни с одной из оцениваемых моделей, а его оценки нуждаются в валидации по русскоязычной человеческой разметке; согласие аннотаторов между собой измеряется по методологии~\cite{artstein2008intercoder} коэффициентом Криппендорфа $\alpha$~\cite{krippendorff1970estimating,krippendorff2011alpha}, восходящим к $\pi$ Скотта~\cite{scott1955reliability}, с порогами интерпретации~\cite{landis1977measurement,carletta1996assessing,krippendorff2004reliability}; для свободных ответов согласие измеряется попарной F-мерой, которая при большом числе отрицательных случаев сходится к $\kappa$~\cite{hripcsak2005agreement}. Часть расхождений аннотаторов~--- законная вариативность, а не шум~\cite{plank2022problem}, что и мотивирует адъюдикацию вместо отбраковки.

\textbf{Временная динамика вопросов и построение выборок.}
Таксономия FreshQA~\cite{vu2023freshllms} делит вопросы на never-changing, slow-changing, fast-changing и вопросы с ложной предпосылкой; эталонная выборка настоящей работы ограничена первым классом и вторым при явной дате-якоре. Обзор темпорального QA~\cite{piryani2025survey} выделяет темпоральность вопроса как самостоятельное измерение, а набор RecencyQA~\cite{piryani2026recency} показывает, что стационарность ответа~--- размечаемое свойство вопроса. Китайские наборы CDQA~\cite{xu2024cdqa} и ChronoQA~\cite{chen2025chronoqa} (5176 вопросов по 300~тыс. китайских новостей) намеренно строятся из времязависимых вопросов, а DRAGOn~\cite{chernogorskii2025dragon} регулярно пересобирает набор по свежим российским новостям, то есть эти работы решают противоположную задачу; вымышленные новостные события NeoQA~\cite{glockner2025neoqa} исключают утечку знаний из предобучения, но неприменимы к оценке реальной системы над реальным корпусом. Отказ на неотвечаемых вопросах остаётся нерешённой задачей~\cite{kirichenko2025abstentionbench}: системы с поиском отказываются от вопросов, на которые могли ответить, и отвечают там, где корпус не содержит ответа~\cite{zhou2025ralm,liu2025crumq}, поэтому вопросы с правильным ответом-отказом и вопросы, ответ на которые заведомо отсутствует в базе~\cite{dorosario2026penalty}, включаются в выборку. Протокол разметки с несколькими независимыми аннотаторами и раздельными метками отвечаемости и ответа заимствован из Natural Questions~\cite{kwiatkowski2019natural}.

\textbf{Китайско- и русскоязычные ресурсы и метрики генерации.}
Метрики точного совпадения и токенного F1 введены в SQuAD~\cite{rajpurkar2016squad}. BERTScore~\cite{zhang2019bertscore} измеряет семантическую близость с помощью контекстных представлений и устойчив к перефразированию и транслитерации имён; сама передача китайских имён средствами русского языка допускает несколько равноправных вариантов~\cite{vashkyavichus2020translation}, что требует алиасов в эталоне. Метрики точности и полноты цитирования предложены в ALCE~\cite{gao2023alce}: утверждение считается подтверждённым, если хотя бы один процитированный фрагмент влечёт его по модели NLI; при этом в многоязычной выдаче модели предпочитают цитировать источники на языке запроса даже в ущерб релевантности~\cite{ki2025nepotism}. Постановку <<вопрос на одном языке~--- свидетельства на другом>> формализует XOR QA~\cite{asai2020xorqa}; для кросс-языкового RAG показано, что смена языка выдачи снижает качество генерации~\cite{liu2025xrag}. Русскоязычные ресурсы оценки~--- REPA~\cite{pugachev2025repa}, POLLUX с открытыми судьями~\cite{martynov2025pollux} и RusHallu-RAG~\cite{sadkovskii2026rushallu}~--- покрывают генерацию и детекцию галлюцинаций на русском, но не вопросно-ответные системы над иноязычным новостным корпусом. В обзоре~\cite{bagrov2026digital} сопоставлены функциональные возможности цифровых инструментов анализа СМИ Северо-Восточной Азии и описана архитектура исследуемой системы, однако количественная оценка ответов в нём отсутствует~--- этот пробел и закрывает настоящая работа.

\textbf{Основы оценки поиска и статистика сравнения систем.}
Парадигма тестовой коллекции~--- фиксированный корпус, запросы и эталонные оценки релевантности~--- восходит к Крэнфилдским экспериментам~\cite{cleverdon1967cranfield}; уже тогда было показано, что расхождения асессоров почти не меняют относительный порядок систем~\cite{lesk1968relevance,voorhees1998variations}, а первый QA-трек TREC ввёл MRR и ручную проверку ответов~\cite{voorhees1999trec8}. Метрика nDCG~\cite{jarvelin2002cumulated} учитывает позицию опорного документа в выдаче. Парные критерии по запросам корректнее сравнения средних~\cite{hull1993using}: используются критерий знаковых рангов Уилкоксона~\cite{wilcoxon1945individual}, парный бутстреп по примерам~\cite{efron1979bootstrap,koehn2004statistical}, поправка Холма на множественные сравнения~\cite{holm1979simple} и порядковая мера размера эффекта $\delta$ Клиффа~\cite{cliff1993dominance}; корреляция Спирмена~\cite{spearman1904proof} служит критерием валидности метрик. Выбор критерия по типу метрики следует протоколу~\cite{dror2018hitchhiker}; оценка LLM как эксперимент над выборкой вопросов требует парных сравнений и планирования размера~\cite{miller2024adding}, а при менее чем нескольких сотнях вопросов интервалы по центральной предельной теореме занижают неопределённость~\cite{bowyer2025position}. Шум несовершенного судьи можно учесть при проверке гипотез, откалибровав его по человеческой подвыборке~\cite{feng2026noisy}.

\section{Основные понятия и обозначения}\label{sec:prelim}

\subsection{Общие обозначения}

В этом разделе вводятся обозначения, используемые далее без изменений (таблица~\ref{tab:notation}).

\begin{table}[h]
  \caption{Основные обозначения}
  \label{tab:notation}
  \centering
  \small
  \begin{tabular}{@{}lp{0.66\textwidth}@{}}
    \toprule
    Символ & Смысл \\
    \midrule
    $t_0$ & момент фиксации корпуса \\
    $D=\{d_j\}_{j=1}^N$ & снимок корпуса; $d_j=(\iota_j,u_j,\tau_j,s_j,x_j^{zh},x_j^{ru})$~--- стабильный идентификатор, URL, дата публикации $\tau_j\le t_0$, источник, текст оригинала, машинный перевод \\
    $\mathcal{I}$ & множество идентификаторов $\iota_j$; фрагмент любого ретривера отображается в $\iota_j$ \\
    $\mathcal{Q}$ & пространство вопросов на русском языке \\
    $\mathcal{A}=\mathcal{A}^+\sqcup\{\varnothing\}$ & пространство ответов; $\varnothing$~--- отказ <<Я не знаю.>> \\
    $2^{\mathcal{I}}$ & пространство цитирований (множеств идентификаторов) \\
    $a^*(q,t)$, $\Inv(q)$ & истинный ответ при состоянии мира в момент $t$; предикат инвариантности ответа во времени \\
    $G=G^+\sqcup G^{\varnothing}$ & эталонная выборка: отвечаемые и требующие отказа вопросы; $|G^+|=n_+$, $|G^{\varnothing}|=n_\varnothing$, $n=n_++n_\varnothing$ \\
    $G_i=(q_i,A^*_i,S^*_i,\mu_i)$ & запись $G^+$: вопрос, множество эталонных ответов (канонический и алиасы), опорные документы $S^*_i\subset\mathcal{I}$, метаданные \\
    $\mathcal{H}$, $\Heval$ & класс гипотез $\{f:\mathcal{Q}\to\mathcal{A}\times2^{\mathcal{I}}\}$ и оцениваемое подмножество \\
    $\mathcal{L}$ & множество открытых моделей провайдера \\
    $\fprod$, $f_0^{(m)}$, $f_\rho^{(m)}$ & промышленная система; модель $m$ без поиска; модель $m$ с ретривером $\rho_{\mathrm{prod}}$ \\
    $M=\{m_k\}$, $M'$ & семейство метрик и валидированное подсемейство \\
    $h^{(r)}_{i,f}$, $\bar h_{i,f}$ & оценка $r$-го аннотатора ответа $f(q_i)$ и её среднее по $r$ \\
    $G_h$, $\mathcal{H}_h$ & подвыборка для человеческой валидации, $|G_h|=n_h$; оцениваемые людьми системы \\
    $Q(f)$, $f^\star$ & композитный показатель и лучшая система \\
    $\rho_S$, $\alpha$ & коэффициент Спирмена; коэффициент согласия Криппендорфа (не путать с уровнем значимости, для которого используется $\alpha=0{,}05$ явно) \\
    \bottomrule
  \end{tabular}
\end{table}

\begin{definition}[статический вопрос]\label{def:static}
Пусть $a^*(q,t)\in\mathcal{A}$~--- истинный ответ на вопрос $q$ при состоянии мира в момент $t$. Вопрос $q$ называется статическим (инвариантным во времени) относительно $t_0$, если
\begin{equation}\label{eq:inv}
  \Inv(q)\iff\exists\,a^*(q)\in\mathcal{A}^+:\ \forall t\ge t_0\quad a^*(q,t)=a^*(q).
\end{equation}
\end{definition}

\begin{definition}[ответ большой языковой модели]\label{def:answer}
Ответом большой языковой модели на вопрос $q$ называется выход $f(q)=(\hat a,\hat S)$ гипотезы $f\in\mathcal{H}$, завершающим шагом которой является генерация текста языковой моделью~--- с поиском по корпусу или без него. Метрика $m_k(\hat a,\hat S;\,G_i)$ определена на ответах; модель или система входит в неё только через $f$.
\end{definition}

\begin{definition}[эталонная выборка]\label{def:golden}
Эталонной выборкой называется $G=G^+\sqcup G^{\varnothing}$, где $G^+=\{(q_i,A^*_i,S^*_i,\mu_i)\}_{i=1}^{n_+}$~--- статические вопросы с непустым множеством эталонных ответов $A^*_i$ (канонический ответ и алиасы, включая транслитерации), опорными документами $|S^*_i|\ge1$ и метаданными $\mu_i$ (тип, категория, источник, дата-якорь, голоса аннотаторов, сложность, разбиение), а $G^{\varnothing}$~--- вопросы, инвариантные по форме, но не подтверждаемые корпусом, для которых правильным ответом является отказ $\varnothing$; $n_\varnothing\approx0{,}1\,n$.
\end{definition}

\subsection{Предположения}

\begin{assumption}[операционализация инвариантности]\label{as:inv}
Предикат~\eqref{eq:inv} непроверяем непосредственно, поэтому вопрос признаётся статическим, если (а)~ответ полностью определяется корпусом, $a^*(q)=g(q,D)$; (б)~существует непустое $S^*_q\subset\mathcal{I}$, документы которого влекут $a^*(q)$; (в)~вопрос содержит явную дату-якорь. Оценка по меткам $R$ аннотаторов $z_r(q)\in\{0,1\}$:
\begin{equation}\label{eq:invhat}
  \widehat{\Inv}(q)=\I\Bigl[\sum_{r=1}^R z_r(q)=R\Bigr].
\end{equation}
\end{assumption}

\begin{assumption}[правила исключения]\label{as:rules}
До разметки инвариантности к вопросу применяются правила $E_1$--$E_5$ (приложение~\ref{app:rules}): $E_1$~--- прогнозы и гипотетика; $E_2$~--- маркеры <<сейчас>>, <<на данный момент>>, <<последний>>, <<текущая ситуация>>; $E_3$~--- величины без даты-якоря (цены, курсы, индексы, действующие должностные лица); $E_4$~--- ответ зависит от момента запроса или политической оценки; $E_5$~--- ложная предпосылка. В терминах~\cite{vu2023freshllms} допускаются вопросы never-changing, вопросы slow-changing~--- только с датой-якорем; fast-changing и false-premise исключаются.
\end{assumption}

\begin{assumption}[независимость судьи]\label{as:judge}
Модель-судья, реализующая метрики $\mcorr$ и $m_J^{(d)}$, не принадлежит множеству $\mathcal{L}\cup\{g_{\theta_{\mathrm{prod}}}\}$; в противном случае используется ансамбль двух судей. Самопредпочтение судьи проверяется отдельно (гипотеза $H_5$).
\end{assumption}

\begin{assumption}[воспроизводимость прогонов]\label{as:repro}
Генерация выполняется при temperature~$0$ и фиксированном seed; поле \texttt{model} каждого ответа и снимок списка моделей провайдера (\texttt{/v1/models}) сохраняются; при неизвестной модели прогон прерывается (fail-fast), что исключает тихую подмену модели провайдером.
\end{assumption}

\subsection{Постановка задачи}\label{sec:problem}

\textbf{Дано.}
\begin{enumerate}
  \item Момент фиксации $t_0$ и снимок корпуса $D=\{d_j\}_{j=1}^N$ с идентификаторами $\mathcal{I}$.
  \item Пространства $\mathcal{Q}$, $\mathcal{A}=\mathcal{A}^+\sqcup\{\varnothing\}$, $2^{\mathcal{I}}$ и предикат $\Inv(q)$~\eqref{eq:inv} с оценкой~\eqref{eq:invhat}.
  \item Класс гипотез и оцениваемое подмножество
  \begin{align}
    \mathcal{H}&=\{f:\mathcal{Q}\to\mathcal{A}\times2^{\mathcal{I}}\},\qquad f(q)=(\hat a,\hat S),\label{eq:H}\\
    \Heval&=\{\fprod\}\cup\{f_0^{(m)}\}_{m\in\mathcal{L}}\cup\{f_\rho^{(m)}\}_{m\in\mathcal{L}},\label{eq:Heval}
  \end{align}
  где $\fprod=g_{\theta_{\mathrm{prod}}}\circ\rho_{\mathrm{prod}}$~--- промышленная система (Qdrant + SearchAPI + генератор), $f_0^{(m)}=(g_m(q),\varnothing)$~--- открытая модель без поиска, $f_\rho^{(m)}=g_m\circ\rho_{\mathrm{prod}}$~--- та же выдача ретривера с другой моделью (контексты кэшируются; абляция <<модель против ретривера>>).
  \item Семейство метрик $M=\{m_k\}$, $m_k:\mathcal{A}\times2^{\mathcal{I}}\times G_i\to[0,1]\cup\{\NA\}$, <<больше~--- лучше>>:
  \begin{itemize}
    \item по эталону (на $G^+$): $\EM=\max_{a\in A^*_i}\I[\norm(\hat a)=\norm(a)]$; токенный F1 после лемматизации, максимум по алиасам; BERTScore-F1 с мультиязычным энкодером; $\mcorr$~--- LLM-судья <<правильность относительно $A^*_i$>>, шкала $0$--$5$, делённая на~$5$; $m_J^{(d)}$, $d\in\{$faithfulness, answer\_relevancy, completeness, citation\_accuracy, temporal\_correctness, allegory\_idiom\_quality, hallucination\_abstention$\}$~--- существующая рубрика как кандидатное подсемейство;
    \item цитирование: $\CP_{\mathrm{id}}=|\hat S\cap S^*_i|/|\hat S|$ ($\NA$ при $\hat S=\varnothing$), $\CR_{\mathrm{id}}=|\hat S\cap S^*_i|/|S^*_i|$ (для систем без поиска равна~$0$), $\CP_{\mathrm{ent}}=\frac{1}{|\hat S|}\sum_{\iota\in\hat S}\I[x_\iota\models\hat a]$ по модели NLI или судье~\cite{gao2023alce};
    \item отказы: $\AR=\frac{1}{n_\varnothing}\sum_{G^{\varnothing}}\I[\hat a=\varnothing]$ и $1-\FAR$, где $\FAR=\frac{1}{n_+}\sum_{G^+}\I[\hat a=\varnothing]$;
    \item поиск (компонент $\rho$): Hit@$k$, Recall@$k$, MRR, nDCG@$k$ относительно $S^*_i$.
  \end{itemize}
  \item Человеческие оценки $h^{(r)}_{i,f}\in\{0,\dots,5\}$ (правильность, обоснованность) на $G_h\subset G^+$, $|G_h|=n_h$, для систем $\mathcal{H}_h\subset\Heval$; $\bar h_{i,f}$~--- среднее по~$r$.
\end{enumerate}

\textbf{Найти.}
\begin{enumerate}
  \item Эталонную выборку $G$, удовлетворяющую ограничениям ниже.
  \item Валидированное подсемейство метрик
  \begin{equation}\label{eq:Mprime}
    M'=\bigl\{m_k\in M:\ \rho_S(m_k,\bar h)\ge\rho_{\min}\ \text{на}\ G_h\times\mathcal{H}_h\bigr\},\qquad \rho_{\min}=0{,}5,
  \end{equation}
  где $\rho_S$ вычисляется с 95\,\%-м бутстреп-интервалом, а потолком служит среднее попарное согласие аннотаторов. Детерминированные метрики без человеческого аналога (Recall@$k$ и т.\,п.) включаются в $M'$ непосредственно.
  \item Для $f\in\Heval$ и $m_k\in M'$~--- средние по вопросам, композитный показатель и лучшую систему:
  \begin{align}
    \bar m_k(f)&=\frac{1}{|I_k(f)|}\sum_{i\in I_k(f)} m_k\bigl(f(q_i);G_i\bigr),\qquad I_k(f)=\{i:\ m_k\ne\NA\},\label{eq:mbar}\\
    Q(f)&=\sum_{k} w_k\,\bar m_k(f),\qquad w_k\ge0,\ \sum_k w_k=1,\label{eq:Q}\\
    f^\star&=\arg\max_{f\in\Heval} Q(f)\label{eq:argmax}
  \end{align}
  с проверкой значимости парных различий; по умолчанию веса $w_k$ равные.
\end{enumerate}

\textbf{Критерий качества и статистика.}
Первичная метрика фиксируется до эксперимента: $\mcorr$, если она проходит валидацию~\eqref{eq:Mprime}, иначе F1; остальные метрики вторичные. Для пары систем $(f,f')$ на общих вопросах вычисляются разности $\delta_i=m_k(f(q_i))-m_k(f'(q_i))$; для $\bar\delta$ строится перцентильный 95\,\%-й интервал парным бутстрепом по вопросам ($B=10^4$); применяется критерий Уилкоксона знаковых рангов с поправкой Холма при уровне значимости $0{,}05$ по всем $f'\ne\fprod$; размер эффекта~--- $\delta$ Клиффа. Разбивка по стратам (категория, источник, тип вопроса, известность события) приводится описательно. Для судьи фиксируются temperature~$0$ и отчёт о собственной воспроизводимости.

\textbf{Ограничения.}
Снимок $D_{t_0}$ зафиксирован (дата, $N$, хеш списка $\mathcal{I}$); $S^*_i\subset\mathcal{I}$ и присутствует в индексе, поэтому покрытие корпуса равно~$1$ по построению; $n\ge150$ (цель 200--300); не менее трёх доменов, двух регионов и охват не менее 12~месяцев; разбиение dev~(30 вопросов)~/ test. Публикуются идентификаторы, URL, даты и короткие фрагменты-свидетельства, но не полные тексты документов.

\section{Метод}\label{sec:method}

Система оценки строится за девять этапов; каждый этап описан так, чтобы его мог воспроизвести другой исследователь по опубликованным конфигурациям.

\begin{enumerate}
  \item \textbf{Отбор документов.} Стратифицированная выборка из пересечения хранилища веб-страниц и индекса Qdrant по стратам <<домен $\times$ месяц $\times$ категория>> с переизбытком $\times2$; фиксируются $t_0$, $N$ и хеш списка $\mathcal{I}$.
  \item \textbf{Черновики вопросов.} Модель $g_{\mathrm{draft}}\notin\mathcal{L}$ порождает по документу кортеж $(q,a,S,\text{evidence},\text{тип},\text{дата})$. Существующие 388 вопросов~\cite{bagrov2026digital} после перекладки ключей \texttt{warc\_id}~$\to$~\texttt{webpage\_id} служат исходным пулом (seed pool), а не эталоном: у них нет эталонных ответов и признака статичности.
  \item \textbf{Независимая разметка.} $R=3$ аннотатора (не менее одного читающего по-китайски) независимо размечают инвариантность по правилам $E_1$--$E_5$, отвечаемость, тип вопроса и правят ответ.
  \item \textbf{Согласие.} По меткам инвариантности, отвечаемости и типа вычисляется коэффициент Криппендорфа $\alpha$ ($\alpha\ge0{,}67$ допустимо, $\alpha\ge0{,}8$ надёжно); по ответам~--- попарный токенный F1 с порогом $0{,}8$.
  \item \textbf{Сборка $G$.} Вопрос включается при единогласии~\eqref{eq:invhat}, при расхождении~--- после адъюдикации четвёртым экспертом. Пилот на 30 вопросах $\to$ правка инструкции $\to$ полный прогон $\to$ заморозка версии \texttt{v1.0} (sha256, журнал изменений).
  \item \textbf{Метрики $M$.} Реализуются метрики по эталону, цитирования, отказов и поиска из раздела~\ref{sec:problem}; существующая семимерная рубрика подключается как кандидатное подсемейство $\{m_J^{(d)}\}$.
  \item \textbf{Валидация $M'$.} На $G_h$ ($n_h=100$) собираются оценки $h^{(r)}_{i,f}$ для $\mathcal{H}_h$; по~\eqref{eq:Mprime} отбираются метрики с $\rho_S\ge0{,}5$; отдельно проверяется самопредпочтение судьи.
  \item \textbf{Прогоны.} Для всех $f\in\Heval$ при temperature~$0$ и фиксированном seed сохраняются сырые генерации, цитирования, поле \texttt{model}, git~sha и снимок \texttt{/v1/models} (предположение~\ref{as:repro}).
  \item \textbf{Статистика.} Вычисляются $\bar m_k(f)$, $Q(f)$~\eqref{eq:Q}, парные разности с бутстреп-интервалами, критерий Уилкоксона с поправкой Холма и $\delta$ Клиффа; результаты сводятся в таблицу вида таблицы~\ref{tab:main}.
\end{enumerate}

\section{Эксперименты}\label{sec:exp}

Для проверки гипотез запланирована эмпирическая часть, параметры которой фиксируются до начала прогонов.

\textbf{Оцениваемые системы.} $\mathcal{L}=\{$DeepSeek-R1-Distill-Qwen-32B, Qwen3.6-27B, Qwen3.8-27B, GigaChat-20B-A3B, GLM-4.1V-9B-Thinking, Qwen-AgentWorld-35B-A3B$\}$~--- шесть открытых моделей провайдера ИВ~РАН (исключены специализированные Hy-MT2, Laguna-XS и Qwen3-Coder). В двух режимах ($f_0^{(m)}$ и $f_\rho^{(m)}$) вместе с $\fprod$ это даёт $|\Heval|=13$ систем.

\textbf{Размеры.} Эталонная выборка $n=200$--$300$ (не менее $150$), $n_\varnothing\approx0{,}1\,n$, $R=3$ аннотатора, подвыборка человеческой валидации $n_h=100$, разбиение dev/test $=30/(n-30)$.

\textbf{Бюджет вычислений.} Разметка: $3\times250\times5$~мин~$\approx21$~ч на аннотатора. Генерации: $13\times250\approx4750$ вызовов (около 40~ч последовательно). Судья: около 25~ч. При сокращённом варианте (защита зимой): $n=150$, 8~систем, $n_h=60$.

\textbf{Гипотезы.}
\begin{itemize}
  \item $H_1$ (основная): $Q(\fprod)>Q(f_0^{(m)})$ для всех $m\in\mathcal{L}$.
  \item $H_2$ (сильная): $Q(\fprod)>Q(f_\rho^{(m)})$ для всех $m\in\mathcal{L}$; при опровержении вклад ретривера и модели разделяется явно.
  \item $H_3$: $Q(f_\rho^{(m)})>Q(f_0^{(m)})$~--- поиск по корпусу повышает качество ответов.
  \item $H_4$: $\rho_S(\mcorr,\bar h)\ge0{,}5$ и $\alpha_{\mathrm{human}}\ge0{,}67$~--- метрика правильности по эталону и разметка надёжны.
  \item $H_5$: самопредпочтение судьи статистически незначимо.
\end{itemize}

\textbf{План работ.} Сентябрь--октябрь 2026~--- постановка, инструкция разметки, пилот на 30 вопросах, обзор литературы; ноябрь--декабрь~--- разметка 250 вопросов и заморозка \texttt{v1.0}; январь 2027~--- прогоны; февраль~--- валидация метрик и статистика; март--апрель~--- текст и слайды; защита~--- июнь 2027.

\begin{table}[h]
  \caption{Заготовка главной таблицы результатов. В ячейках~--- $\bar m_k(f)$ с 95\,\%-м бутстреп-интервалом; звёздочкой отмечаются различия с $\fprod$, значимые после поправки Холма ($p_{\mathrm{Holm}}<0{,}05$). Значения будут заполнены после прогонов.}
  \label{tab:main}
  \centering
  \small
  \begin{tabular}{@{}llcccccc@{}}
    \toprule
    Система & Режим & $\mcorr$ & F1 & $\CP_{\mathrm{id}}$ & $\AR$ & $1-\FAR$ & $Q(f)$ (ранг) \\
    \midrule
    $\fprod$ & Qdrant + SearchAPI & $\dots$ & $\dots$ & $\dots$ & $\dots$ & $\dots$ & $\dots$ \\
    \midrule
    $f_0^{(m_1)}$ & без поиска & $\dots$ & $\dots$ & $\NA$ & $\dots$ & $\dots$ & $\dots$ \\
    $f_0^{(m_2)}$ & без поиска & $\dots$ & $\dots$ & $\NA$ & $\dots$ & $\dots$ & $\dots$ \\
    $\vdots$ & $\vdots$ & $\vdots$ & $\vdots$ & $\vdots$ & $\vdots$ & $\vdots$ & $\vdots$ \\
    \midrule
    $f_\rho^{(m_1)}$ & $\rho_{\mathrm{prod}}$ & $\dots$ & $\dots$ & $\dots$ & $\dots$ & $\dots$ & $\dots$ \\
    $f_\rho^{(m_2)}$ & $\rho_{\mathrm{prod}}$ & $\dots$ & $\dots$ & $\dots$ & $\dots$ & $\dots$ & $\dots$ \\
    $\vdots$ & $\vdots$ & $\vdots$ & $\vdots$ & $\vdots$ & $\vdots$ & $\vdots$ & $\vdots$ \\
    \bottomrule
  \end{tabular}
\end{table}

\section{Обсуждение}\label{sec:disc}

\textbf{Самопредпочтение судьи.} Фактический генератор промышленной системы~--- DeepSeek-R1-Distill-Qwen-32B, входящая и в $\mathcal{L}$. Если судья окажется той же или родственной моделью, оценки $\mcorr$ будут смещены в пользу её ответов~\cite{panickssery2024selfpreference,li2025preference}. Меры: предположение~\ref{as:judge}, гипотеза $H_5$, панель судей разных семейств~\cite{verga2024replacing} и опора на детерминированные метрики $\EM$, F1, $\CP_{\mathrm{id}}$ как контроль.

\textbf{Подмена модели провайдером.} Локальный провайдер (llama-swap) может тихо заменить неизвестную модель на модель по умолчанию, что уже наблюдалось для конфигурации по умолчанию. Мера: fail-fast, запись поля \texttt{model} каждого ответа и снимка списка моделей (предположение~\ref{as:repro}).

\textbf{Утечка знаний в предобучение.} События 2022--2023 годов могли попасть в обучающие данные открытых моделей, и высокий результат $f_0^{(m)}$ может отражать память модели, а не рассуждение по корпусу; известно, что на популярных сущностях параметрическая память конкурентоспособна с поиском~\cite{mallen2022popqa}. Мера: страта <<известность события>> в метаданных $\mu_i$ и описательная разбивка результатов по ней; вопросы $G^{\varnothing}$ дополнительно проверяют склонность к уверенному ответу без опоры~\cite{dorosario2026penalty}.

\textbf{Транслитерация имён и терминов.} Русская передача китайских имён и названий неоднозначна~\cite{vashkyavichus2020translation}, что ломает $\EM$ и токенный F1. Мера: алиасы в $A^*_i$ и мягкие метрики (BERTScore, $\mcorr$); согласие судьи с человеком на таких вопросах проверяется отдельно.

\textbf{Область применимости.} В названии работы указаны страны Восточной Азии, тогда как корпус исследуемой системы на текущем этапе китайскоязычный. Методика не зависит от языка корпуса, а расширение выборки на японские и корейские источники требует лишь повторения этапов~1--5 раздела~\ref{sec:method} с аннотаторами, владеющими соответствующими языками. Метрики поиска воспроизводимы только при доступе к индексу ИВ~РАН; все прочие метрики воспроизводимы по опубликованным данным.

\section{Заключение}\label{sec:concl}

В работе поставлена задача оценки качества ответов больших языковых моделей на статические вопросы по новостному корпусу стран Восточной Азии: введён предикат инвариантности $\Inv(q)$, определены эталонная выборка $G=G^+\sqcup G^{\varnothing}$, класс гипотез $\Heval$, семейство метрик $M$ и композитный показатель $Q(f)$ с процедурой статистического сравнения. Предложена девятиэтапная методика построения системы оценки, включающая независимую разметку тремя аннотаторами, измерение согласия и валидацию метрик по человеческим оценкам. Сформулированы гипотезы $H_1$--$H_5$ и план экспериментов над промышленной системой~\cite{bagrov2026digital} и шестью открытыми моделями в двух режимах. Результатом работы станут опубликованная эталонная выборка, валидированное подсемейство метрик $M'$ и главная таблица результатов с интервалами и поправкой на множественные сравнения.

%%%%%%%%%%%%%%%%%%%%%%%%%%%%%%%%%%%%%%%%%%%%%%%%%%%%%%%%%%%%

\bibliographystyle{unsrtnat}
\bibliography{references}

%%%%%%%%%%%%%%%%%%%%%%%%%%%%%%%%%%%%%%%%%%%%%%%%%%%%%%%%%%%%

\newpage
\appendix
\section{Приложение}\label{app}

\subsection{Правила исключения вопросов $E_1$--$E_5$}\label{app:rules}

Правила применяются аннотаторами до разметки инвариантности; вопрос, к которому применимо хотя бы одно правило, в $G^+$ не включается.

\begin{enumerate}
  \item[$E_1$] \textbf{Прогнозы и гипотетика.} Вопросы о будущем и условные конструкции (<<что произойдёт, если~\dots>>, <<как изменится~\dots>>).
  \item[$E_2$] \textbf{Маркеры текущего момента.} Слова и обороты <<сейчас>>, <<на данный момент>>, <<последний>>, <<новейший>>, <<текущая ситуация>> и их эквиваленты, привязывающие ответ к моменту запроса.
  \item[$E_3$] \textbf{Величины без даты-якоря.} Цены, курсы валют, биржевые индексы, численность, действующие должностные лица без указания даты, на которую запрашивается значение. Вопрос сохраняется, если дата-якорь добавлена явно (<<по данным на 12~мая 2023~года>>).
  \item[$E_4$] \textbf{Зависимость от момента запроса или оценки.} Ответ требует политической, моральной или иной оценочной позиции либо меняется в зависимости от того, когда задан вопрос.
  \item[$E_5$] \textbf{Ложная предпосылка.} Вопрос предполагает событие или факт, которого не было; такие вопросы могут попасть в $G^{\varnothing}$ только при явной пометке типа.
\end{enumerate}

В терминах таксономии~\cite{vu2023freshllms}: допускаются вопросы never-changing; slow-changing~--- только с датой-якорем; fast-changing и false-premise исключаются.

\subsection{Схема записи эталонной выборки}\label{app:schema}

Записи $G$ хранятся в формате JSON Lines; одна строка~--- одна запись. Значения строковых полей на русском и китайском языках в примере заменены многоточием.

\begin{lstlisting}[language=json]
{
  "id": "ea-000001",
  "version": "1.0",
  "split": "test",
  "question": {"ru": "...", "zh": "...", "en": "..."},
  "question_type": "factoid",
  "category": "economy",
  "answerable": true,
  "answer": {
    "canonical": "...",
    "aliases": ["...", "..."],
    "type": "entity",
    "lang": "ru"
  },
  "supporting_docs": [
    {
      "webpage_id": "...",
      "urn_uuid": "urn:uuid:...",
      "uri": "https://...",
      "published_at": "2023-05-12",
      "source": "thenewslens.com",
      "evidence_span_ru": "...",
      "evidence_span_zh": "..."
    }
  ],
  "requires_multi_doc": false,
  "time_invariance": {
    "label": "never-changing",
    "anchor_date": "2023-05-12",
    "votes": [1, 1, 1],
    "exclusion_rules_checked": ["E1", "E2", "E3", "E4", "E5"]
  },
  "difficulty": "medium",
  "provenance": {
    "draft_model": "...",
    "draft_prompt_id": "draft-v1",
    "seed_pool_ref": null,
    "annotators": ["a1", "a2", "a3"],
    "adjudicated": false,
    "created_at": "2026-11-01T00:00:00Z"
  }
}
\end{lstlisting}

Поле \texttt{answerable} равно \texttt{false} для записей $G^{\varnothing}$; в них \texttt{answer.canonical} принимает значение отказа, а \texttt{supporting\_docs} пусто. Поле \texttt{time\_invariance.votes} содержит метки $z_r(q)$ трёх аннотаторов, по которым вычисляется~\eqref{eq:invhat}.

\end{document}
